\originalchapter{核反应、中子与核能}\label{ch:reactions}
\originalsection{反应阈值与截面}
吸能反应 $a+A\to b+B$ 的精确实验室阈值, 由末态在质心系相对动量为零取得. 用 $s=m_a^2+m_A^2+2m_AE_a$ ($c=1$), 得
\begin{equation}
 T_{a,\rm th}=\frac{(m_b+m_B)^2-(m_a+m_A)^2}{2m_A}
 \simeq -Q\left(1+\frac{m_a}{m_A}\right).
\end{equation}
右边近似要求 $|Q|$ 相对核质量小. 阈值之外的库仑势垒会让截面仍然很小, 二者不应混淆.

设入射通量为 $\Phi$、靶核数密度为 $n$, 厚度 $\dd x$ 内反应比例为 $n\sigma\dd x$. 不受再生粒子影响的窄束满足 \begin{equation}\Phi(x)=\Phi(0)\e^{-n\sigma x}.\end{equation}平均自由程 $\ell=(n\sigma)^{-1}$. 混合材料用宏观截面 $\Sigma=\sum_i n_i\sigma_i$. $1$ barn 为 $10^{-28}\,\mathrm m^2$, 换算时不可同 fm$^2$ 混用.

\originalsection{中子慢化}
质量为 $m$ 的中子与静止质量 $Am$ 的核弹性碰撞. 将质心平移速度 $v_0/(A+1)$ 与质心系中子速度 $Av_0/(A+1)$ 相加, 得实验室末动能比
\begin{equation}
 \frac{E_1}{E_0}=\frac{A^2+1+2A\cos\theta}{(A+1)^2},
 \qquad \alpha\le E_1/E_0\le1,\quad\alpha=\left(\frac{A-1}{A+1}\right)^2.
\end{equation}
$\theta$ 是质心散射角. 对氢 $A=1$, 一次碰撞可转移几乎全部动能; 对重核比例很小. 若质心散射各向同性, $E_1/E_0$ 在 $[\alpha,1]$ 上均匀, 平均对数能损
\begin{equation}
 \xi=\left\langle\log\frac{E_0}{E_1}\right\rangle
 =1+\frac{\alpha\log\alpha}{1-\alpha}.
\end{equation}
从 $E_0$ 慢化到 $E$ 的平均碰撞次数约为 $\log(E_0/E)/\xi$. 好慢化剂既要质量轻以有效减速, 又要吸收截面小以保留中子; 因而不能只凭原子质量排列优劣.

\originalsection{裂变与链式反应}
重核裂成两个中质量核, 每核子的结合更紧, 每次裂变总释放能量约 $200\,\mathrm{MeV}$. 大部分先表现为裂变碎片动能, 另有 prompt 中子、光子和随后 $\beta$ 衰变能. 裂变碎片通常中子过多, 所以放射性强; 迟发中子来自少量碎片的后续衰变.

有效增殖因子 $k_{\rm eff}$ 是相邻中子世代数之比. $k<1$ 为次临界, $k=1$ 为临界, $k>1$ 为超临界. 这个定义还包含泄漏、非裂变吸收与能谱变化; 裂变每次平均产生多于一个中子不保证整个装置超临界. 迟发中子把部分动力学时间尺度拉长, 是反应堆控制的重要原因.

核能利用讨论在本课程限于能量与反应动力学, 不提供装置设计或运行操作程序. 活度与吸收剂量是不同量: 活度是每秒衰变次数, 吸收剂量是每单位质量吸收的能量. 测得一条光子线还需探测效率和几何接受度才能推回源活度.

\originalsection{聚变与热平均}
常用入门反应为 D+T$\to{}^4$He+n ($17.6\,\mathrm{MeV}$), D+D$\to{}^3$He+n ($3.27\,\mathrm{MeV}$) 或 T+p ($4.03\,\mathrm{MeV}$), D+$^3$He$\to{}^4$He+p ($18.35\,\mathrm{MeV}$). 在通常受控热核能区, D--T 截面较大且适合较低温度, 但这些是反应性质, 不是对今日技术进度的判断.

热反应率密度是 $n_1n_2\langle\sigma v\rangle$; 同种反应物须再除 $2$, 避免每对重复计数. 库仑隧穿因子近似 $\e^{-b/\sqrt E}$, Maxwell 因子为 $\e^{-E/T}$, 两者使贡献集中在 Gamow 能区. 令指数 $E/T+b/\sqrt E$ 对 $E$ 的导数为零, 得 $E_G=(bT/2)^{2/3}$. 最有效能量一般高于简单的平均热能, 不能以 $\sigma(T)$ 代替整个热平均.

太阳 pp 链先经弱过程 p+p$\to$d+$e^++\nu_e$, 再经 p+d$\to{}^3$He+$\gamma$. 常见主链以两份 ${}^3$He 合成 ${}^4$He 与两个质子; 较稀有 hep 支路 p+$^3$He$\to{}^4$He+$e^++\nu_e$. 原题指定哪条支路就保留哪条, 不能为了写成常见 pp 主链而改掉题目.
